\section{Experimental Setup}
\label{sec:experimental_setup}

\subsection{Data and chronological evaluation}
Table~\ref{tab:datasets} lists the evaluated multivariate series. The main set contains ETTh1, ETTh2, ETTm1, and ETTm2 from the Electricity Transformer Temperature repository~\cite{ettdata}, the Appliances Energy Prediction dataset~\cite{appliancesdata}, and a subset of Building Data Genome 2 (BDG2)~\cite{bdg2data}. We also evaluate two additional BDG2 site subsets, Fox and Panther. The ETT series contain seven channels. For Appliances, we remove the timestamp and the two random variables \texttt{rv1} and \texttt{rv2}, leaving 26 numerical channels. Forecasts and metrics cover all retained channels.

The BDG2 subsets each contain 15 electricity meters. The main subset uses site Rat, denoted \texttt{bdg2} in the figures and tables. Within each site, meters were selected by the lowest missing-value fractions over the supplied series. The released configuration fixes the meter IDs and order. Fox and Panther are reproduced by forward filling and then backward filling the pinned cleaned source. The archived Rat input differs from that procedure at 78 cells; the release supplies explicit snapshot adjustments and verifies the resulting numerical hash. The original cleaning decisions behind those differences could not be recovered. These adjustments preserve the evaluated input rather than define a new cleaning rule. Meter selection uses whole-series missingness, and backward filling can use later observations. We therefore distinguish causal model updates on these prepared series from a strictly causal raw-data acquisition pipeline.

All datasets use the same lookback length $L=96$, forecast horizon $H=24$, and forecast-origin stride of 24 time steps. These are sample counts rather than a common duration. ETTh and BDG2 are hourly, ETTm is sampled every 15 minutes, and Appliances every 10 minutes. A complete 24-step target block is scored before the next block is issued.

For a series of length $T$, let $n=\lfloor T/2\rfloor$ and $a=\lfloor0.8n\rfloor$. The evaluation targets start at index $n$. Their origins are $n,n+24,\ldots$, retaining only complete target blocks. Channel means and standard deviations are computed from indices $[0,n)$ and held fixed thereafter. Normalization divides by the standard deviation plus $10^{-8}$. Thus, the normalization interval includes the validation interval $[a,n)$, although it excludes evaluation targets. We report errors on this normalized scale. These chronological splits are the splits used in this study, rather than the standard ETT benchmark partitions.

\begin{table}[!t]
\caption{Prepared series used for evaluation. Evaluation start is a zero-based sample index. The Blocks column gives the number of complete target blocks. Fox and Panther are additional sites from the same BDG2 corpus. The common lookback, horizon, and stride are specified in the text.}
\label{tab:datasets}
\centering
% Dataset-specific quantities. Common L, H, and stride are defined in Section IV.
\begingroup\fontsize{9}{11}\selectfont
\setlength{\tabcolsep}{2.3pt}
\begin{tabular}{lrrrr}
\toprule
Series & Samples & $C$ & Eval. & Blocks \\
\midrule
ETTh1 & 17420 & 7 & 8710 & 362 \\
ETTh2 & 17420 & 7 & 8710 & 362 \\
ETTm1 & 69680 & 7 & 34840 & 1451 \\
ETTm2 & 69680 & 7 & 34840 & 1451 \\
Appliances & 19735 & 26 & 9867 & 411 \\
BDG2 (Rat) & 17544 & 15 & 8772 & 365 \\
BDG2 (Fox) & 17544 & 15 & 8772 & 365 \\
BDG2 (Panther) & 17544 & 15 & 8772 & 365 \\
\bottomrule
\end{tabular}
\endgroup

\end{table}

\subsection{Base models and evaluation groups}
We use DLinear~\cite{dlinear2022} and a compact channel-independent PatchTST implementation~\cite{patchtst2023}. DLinear uses a moving-average kernel of 25 with replicate padding and separate linear maps for the trend and residual components. Its temporal maps are shared across channels. PatchTST uses patches of length 16 with stride 8, an embedding dimension of 64, four attention heads, and two encoder layers. The feed-forward dimension is 128 and dropout is zero. A learned position embedding and a linear forecast head complete the network. The implementation also includes a learned affine input transformation per channel. These architectures define our comparison setting and are not presented as exact reproductions of all configurations in the original papers.

The existing base-model evaluation contains two training variants. In the \emph{legacy} variant, the model is trained on $[0,a)$ with Adam, a learning rate of $10^{-3}$, batch size 32, and mean squared error as the loss. Training windows are randomly sampled. The number of updates is selected by validation mean squared error from $\{200,500,1000,2000,4000,8000,20000\}$. In the \emph{refit} variant, the same architecture is initialized again with the corresponding seed and trained on $[0,n)$ for that selected number of updates. Both variants remain fixed during evaluation. Six datasets, two architectures, three seeds ($0,1,2$), and two variants give 72 base-model conditions. Fox and Panther add four conditions using the legacy variant and seed 0. These site results are reported separately.

A second group contains newly trained models for the learned-adapter and feature-selection comparisons. It uses the six main datasets, the two architectures, and the same three seed values. Training and validation samples each contain a previous context--target pair and a current context--target pair. Both pairs lie entirely within the relevant interval, $[0,a)$ or $[a,n)$. Within an interval starting at $u$, current forecast origins begin at $u+L+H$ and advance by $H$. During evaluation, the first origin remains $n$, and historical observations before $n$ provide its input context. This group is kept separate from the legacy/refit group because its training objective and budget differ.

\subsection{Comparators and controlled training}
The uncorrected base forecast is the reference for total correction benefit. The proposed method uses the fixed configuration in Section~\ref{sec:method}, including $K=4$. For each condition, comparison methods receive the same base forecasts and evaluation targets. Online regression statistics, residual histories, and blending statistics are initialized at the evaluation boundary. Offline-trained adapter weights are retained. Initial blocks are included in the main metrics.

\subsubsection{Residual representation controls}
The coefficient-only control removes the endpoint input while retaining the four DCT components, current forecast coefficients, ridge update, and blending procedure. Three additional controls replace the endpoint with the first residual, the twelfth residual, or the mean residual of the preceding block. The twelfth residual is the selected middle position for the even horizon length. All four scalar-summary variants use the same $\sqrt H$ scaling, regression dimensions, and retained array sizes. For the block mean $m_{t-1,c}$, $\sqrt H\,m_{t-1,c}=z_{t-1,0,c}$. Thus, the mean is recoverable from the complete retained coefficient vector. However, each component's regression receives only its own residual coefficient $z_{t-1,j,c}$. The mean control duplicates that input only for $j=0$; for $j>0$, it additionally supplies the DC coefficient to a non-DC regression. The feature comparisons use fixed forecasts from the jointly trained width-64 models, without additional training or checkpoint selection.

\subsubsection{Learned residual adapters}
We implement the two-projection residual adapter described by TEFL~\cite{tefl2026}, with a rectified linear activation, no projection biases, and shared weights across channels. We evaluate hidden widths 2, 4, and 64. Width 64 is an expanded hidden representation at $H=24$, rather than a rank reduction. The second projection is initialized to zero.

Training begins with three chronological warmup epochs of the base model. One condition uses mean absolute error (MAE) alone. The other adds spectral flatness (SF) with coefficient one. For SF, consecutive residual blocks within each batch are concatenated along time. We compute full Fourier power plus $10^{-8}$, take its geometric-to-arithmetic mean ratio over time frequencies, and average the ratio over channels. The constant frequency is retained, and the final short batch is included. This specifies our multivariate reduction of the SF term.

The second stage uses 12 epochs of MAE optimization, either for the base model alone or jointly for the base and adapter. Both stages use AdamW with learning rate $10^{-3}$, weight decay 0.01, and batch size 32, without a scheduler. The second stage shuffles training samples. Paired runs share the base initialization, warmup state, and shuffle seed. Gradients pass through the base forecasts used to construct training residuals. Checkpoints are selected by validation MAE over epochs 0--12, retaining the earlier epoch on a tie. The first validation forecast uses a zero previous residual, matching evaluation initialization. The final three-epoch warmup state is shared across the paired second-stage runs.

At evaluation, adapter weights are fixed and the previous observed base residual is supplied as input. We report its raw correction and its correction with the same global blending controller used by the proposed method. The principal correction comparison applies both methods to the same selected joint base model. A separate comparison with the corresponding base-only training run checks the effect of joint training. Training samples and update counts are matched within each pair, but computation per update is not assumed equal. The asymmetric context and horizon, compact backbones, and explicitly defined SF reduction make this a transfer of the TEFL training design rather than a complete reproduction.

\subsubsection{Online Fourier regression}
The ELF comparisons use cumulative ridge regression with coefficient 20 and one newly completed training window per forecast block. Regression begins without pre-evaluation training pairs. Each input window is centered by its channel mean, which is also subtracted from its target window. Fourier coefficients use orthonormal scaling, and the input mean is restored after reconstruction. Before the first fitted regression, the auxiliary forecast repeats the most recent daily cycle. Channel scales are updated from the initial context and subsequently observed values. In the unscaled sufficient statistics, the ridge term for channel $c$ is $20\sigma_c^2 I$, where $\sigma_c$ is the current channel scale. Linear systems are solved directly.

We evaluate seven input/output frequency counts $(d,q)$: $(87,10)$, $(3,10)$, $(5,10)$, $(9,10)$, $(87,4)$, $(3,4)$, and $(5,4)$. Here, $d$ counts signed input frequencies including the constant term, and $q$ counts retained real-FFT output frequencies including the constant term. Their labels are, respectively, \texttt{full}, \texttt{in3}, \texttt{in5}, \texttt{in9}, \texttt{out4}, \texttt{compact3}, and \texttt{compact5}. The principal comparison uses the shared global controller. Auxiliary evaluations also use the transferred ELF weighting rule with fast, slow, and merging weights, learning rate 0.5, and a five-update fast window. That rule issues the base forecast for the first five blocks. Its loss scales use past daily differences, with seasonal periods 24, 96, and 144 for hourly, 15-minute, and 10-minute data, respectively.

For storage accounting, conjugate input frequencies are converted to independent real coordinates and the zero centered constant input is removed. We store the upper triangle of the real symmetric Gram matrix and real target cross-statistics. This changes representation and storage, not the intended ridge prediction. Both this regression and the proposed regression use 64-bit arithmetic. The learned adapters use 32-bit weights. We count their actual retained arrays under the exclusions in Section~\ref{sec:method}. The ELF configurations adapt the source method~\cite{elf2025} to our window lengths, update schedule, and base models.

\subsection{Input ablations and retrospective selection sensitivity}
We additionally replay 148 fixed forecast streams: the 72 legacy/refit conditions, four additional-site conditions, and 36 width-64 joint conditions for each of the MAE and SF warmups. Input ablations retain only the endpoint, remove the residual DCT coefficients, or remove the current forecast coefficients. Together with the complete and coefficient-only methods, these regressions share the four-component output subspace, regularized intercept, and update/controller settings. Their input dimensions and retained sizes differ. An endpoint-persistence control repeats the preceding endpoint across the horizon without a fitted regression. Both raw and globally blended predictions are scored.

The setting sensitivity contains 13 candidates: the reference and one-factor changes to $K\in\{1,2,3,6,8\}$, ridge coefficient $\{0.1,10\}$, half-life $\{32,512,\infty\}$ blocks, or controller window $\{8,128\}$ blocks. A second pool contains all 21 candidates, additionally including the three scalar-summary controls and five input/persistence controls. This is not an exhaustive interaction search or a memory-matched comparison.

For a stream of $N$ blocks, the principal split selects the candidate with the smallest mean global MSE on $[0,\lfloor N/2\rfloor)$ and scores it on $[\lfloor N/2\rfloor,N)$. Additional expanding-prefix selections score the second, third, and fourth quarters separately, using every preceding block for selection. Each candidate's regression and controller state continues across the boundary, with prediction before target update; the chosen candidate stays fixed within the scored interval. We compare the fixed reference, prefix-selected candidate, and a retrospective oracle selected using the scored interval itself. The oracle gap describes only the declared finite pool. These already-used evaluation series do not become independent validation through retrospective splitting. Appendix~\ref{app:sensitivity} documents the inherited choices and full records.

To assess dependence on the Rat snapshot, we also exclude all \texttt{bdg2} conditions from the full-period input and summary comparisons. The remaining five datasets give 60 legacy/refit conditions and ten seed-averaged dataset--backbone pairs per joint-base warmup. Forecasts, model settings, and all other inputs remain unchanged. This is a retrospective dataset-exclusion sensitivity analysis, not a re-evaluation of a different Rat preprocessing rule or a new independent test. Fox and Panther are separate additional-site results and are not included in these main-group aggregates.

\subsection{Metrics, aggregation, and reproducibility}
For $N$ evaluated blocks, define $E_{t,h,c}=\widehat y_{t,h,c}-y_{t,h,c}$. We compute
\begin{equation}
 \begin{split}
 \operatorname{MSE}&=\frac{1}{NHC}\sum_{t,h,c}E_{t,h,c}^{2},\\
 \operatorname{MAE}&=\frac{1}{NHC}\sum_{t,h,c}|E_{t,h,c}|.
 \end{split}
 \label{eq:eval_metrics}
\end{equation}
For a specified comparator with loss $\ell_{\mathrm{ref}}$, the relative reduction is $100(1-\ell_{\mathrm{prop}}/\ell_{\mathrm{ref}})$. Positive values favor the proposed method. MSE and MAE reductions are computed separately.

The overall legacy/refit accuracy--storage summary averages reductions over its 72 conditions. Condition-wise plots first average seed losses within each dataset, architecture, and training variant before computing reductions. The learned-adapter and feature comparisons also average the three seed losses first, giving 12 paired dataset--architecture conditions for each training configuration. We report their mean reductions and paired win counts. Consequently, an average of condition-wise reductions need not equal a reduction calculated from pooled losses. Seeds, architectures, training variants, and BDG2 sites share data and are not treated as independent datasets. We do not infer statistical significance from a count of wins.

For the local failure illustration, we examine all contiguous 32-block windows starting at block index 32 or later, using zero-based indices. We select the largest relative MSE increase over the uncorrected base across all 76 legacy/refit and additional-site conditions. This is a deliberately selected worst interval, not a representative random segment. Full-period results retain the initial blocks and all adverse conditions.

Saved predictions are rescored to reproduce the reported losses. Checkpoint replay checks the selected validation scores and base forecasts. Additional checks cover prediction-before-update ordering, numerical equivalence of the packed regression, and equivalence of the real and complex ELF representations. Source and input hashes identify the evaluated versions. These checks establish computational consistency within the experiment. The evaluation series were also used during method development and candidate selection, so the results are not described as an untouched external confirmation. We report this distinction alongside the chronological separation of fitting and evaluation targets.
